% !TEX root = ../../../main.tex
\section{Maximum (extreme value theorem)}
\label{sec:analysis-i:extreme-value-theorem}

% BEGIN TUFTE PLACEHOLDER: supplied sample, lines 743--764.
\begingroup\sloppy
% Replace this entire specimen block with reviewed course notes.
\newthought{Occasionally} \LaTeX{} will generate an error message:\label{template:err:too-many-floats}
\begin{docspec}
  Error: Too many unprocessed floats
\end{docspec}
\LaTeX{} tries to place floats in the best position on the page.  Until it's
finished composing the page, however, it won't know where those positions are.
If you have a lot of floats on a page (including sidenotes, margin notes,
figures, tables, etc.), \LaTeX{} may run out of ``slots'' to keep track of them
and will generate the above error.

\LaTeX{} initially allocates 18 slots for storing floats.  To work around this
limitation, the \TL document classes provide a \doccmddef{morefloats} command
that will reserve more slots.

The first time \doccmd{morefloats} is called, it allocates an additional 34
slots.  The second time \doccmd{morefloats} is called, it allocates another 26
slots.

The \doccmd{morefloats} command may only be used two times.  Calling it a
third time will generate an error message.  (This is because we can't safely
allocate many more floats or \LaTeX{} will run out of memory.)

\lipsum[6]
\par\endgroup
% END TUFTE PLACEHOLDER
